\chapter{Tổng quan}
\label{Chapter2}

\section{Tổng quan về hệ thống thuê nhà trọ trực tuyến}

Các hệ thống tìm kiếm và giao dịch bất động sản cho thuê trực tuyến giữ vai trò quan trọng trong việc kết nối cung và cầu, trong đó phân khúc nhà trọ và phòng cho thuê bình dân chiếm tỷ trọng lớn tại các nước đang phát triển, đặc biệt là Việt Nam. Phần này trình bày khái quát các khái niệm cơ bản liên quan đến thị trường cho thuê nhà trọ, cùng các mô hình hệ thống công nghệ đang vận hành phổ biến hiện nay.

\FloatBarrier
\subsection{Khái niệm nhà trọ và thị trường cho thuê nhà trọ}

Nhà trọ hay phòng trọ cho thuê là một loại hình bất động sản cư trú quy mô nhỏ, được xây dựng hoặc phân tách từ một ngôi nhà nguyên căn nhằm mục đích cho cá nhân hoặc hộ gia đình thuê lại để ở trong một thời hạn nhất định. Đối tượng thuê nhà trọ chủ yếu tập trung vào các nhóm dân cư có tính dịch chuyển cao và thu nhập trung bình thấp như: sinh viên các trường đại học, người lao động ngoại tỉnh nhập cư, công nhân tại các khu công nghiệp, và các hộ gia đình trẻ chưa có đủ tích lũy tài chính để sở hữu nhà riêng.

Thị trường cho thuê nhà trọ tại các đô thị lớn của Việt Nam (đặc biệt là Hà Nội và TP. Hồ Chí Minh) mang tính đặc thù cao:
\begin{itemize}
    \item \textbf{Cầu vượt cung và mất cân bằng thông tin:} Lượng người đổ về các thành phố lớn học tập và làm việc tăng trưởng liên tục qua các năm, trong khi quỹ đất và nguồn cung phòng trọ giá rẻ có giới hạn. Điều này tạo nên thị trường do người bán quyết định, dẫn tới tình trạng chủ nhà trọ nắm thế chủ động và người thuê dễ rơi vào thế bị động khi tìm phòng.
    \item \textbf{Sự phân tán dữ liệu:} Thông tin phòng trọ thường không được quản lý tập trung mà phân tán nhỏ lẻ qua các biển treo thủ công, các hội nhóm mạng xã hội (Facebook, Zalo) hoặc các trang rao vặt không chuyên.
    \item \textbf{Sự thiếu minh bạch về chi phí và chất lượng:} Giá thuê phòng, giá điện, giá nước và các phí dịch vụ kèm theo (vệ sinh, bảo vệ, internet) thường do chủ nhà tự quyết định và không có hệ thống đối chiếu rõ ràng. Chất lượng phòng trọ trên ảnh quảng cáo thường sai lệch rất nhiều so với thực tế bàn giao.
\end{itemize}

\FloatBarrier
\subsection{Các mô hình hệ thống thuê nhà hiện nay}

Hiện nay, các giải pháp tìm kiếm và giao dịch bất động sản cho thuê trực tuyến tại Việt Nam và trên thế giới đang vận hành theo ba mô hình công nghệ chính:
\begin{enumerate}
    \item \textbf{Mô hình bảng tin rao vặt truyền thống:} Các website vận hành như một bảng tin công cộng (ví dụ: Phongtro123, các group Facebook). Người cho thuê đăng tin tự do và người thuê tự liên hệ trực tiếp. Ưu điểm là lượng tin phong phú và chi phí đăng thấp. Nhược điểm lớn là không kiểm soát được chất lượng tin đăng, tỷ lệ tin ảo, tin trùng lặp hoặc lừa đảo cao, gây lãng phí lớn về thời gian cho người dùng.
    \item \textbf{Mô hình sàn giao dịch tích hợp dịch vụ:} Các nền tảng đứng ra làm trung gian kiểm soát thông tin bài đăng và hỗ trợ các tính năng giá trị gia tăng như bản đồ tìm kiếm, thanh toán hóa đơn trực tuyến, và phân quyền môi giới chuyên nghiệp (ví dụ: Nhà Tốt, Batdongsan.com.vn). Mô hình này cải thiện tính minh bạch nhưng phần lớn quy trình vận hành vẫn chạy bằng nhân lực thủ công (nhân viên duyệt tin, nhân viên tư vấn gọi điện), dẫn đến tốc độ xử lý chậm và khó mở rộng quy mô một cách tối ưu.
    \item \textbf{Mô hình nền tảng thông minh tích hợp trí tuệ nhân tạo:} Là mô hình thế hệ mới mà đề tài này hướng đến xây dựng (Thuê Nhà Trọ). Hệ thống tận dụng sức mạnh của các mô hình ngôn ngữ lớn và học máy để tự động hóa hoàn toàn các tác vụ phức tạp: AI Agent hỗ trợ tư vấn và so sánh phòng bằng ngôn ngữ tự nhiên, AI tự động kiểm duyệt hình ảnh và nội dung tin đăng để lọc tin rác, AI định giá dựa trên ngữ cảnh thị trường, và gợi ý tin cá nhân hóa theo hành vi thời gian thực. Mô hình này giúp cải thiện trải nghiệm người dùng, giảm đáng kể chi phí vận hành thủ công và nâng cao độ an toàn giao dịch.
\end{enumerate}

\section{Khảo sát các nền tảng hiện có}

Trước khi đề xuất giải pháp, nhóm tiến hành khảo sát một số nền tảng tiêu biểu trong lĩnh vực rao vặt và tìm kiếm bất động sản cho thuê, bao gồm ba nền tảng phổ biến tại Việt Nam (Batdongsan.com.vn, Phongtro123, Nhà Tốt) và một nền tảng quốc tế đại diện cho hướng ứng dụng dữ liệu và công nghệ ở mức cao (Zillow). Việc khảo sát nhằm nhận diện các tính năng đã được thị trường chấp nhận, đồng thời chỉ ra những hạn chế còn tồn tại để làm cơ sở cho định hướng của đề tài.

\FloatBarrier
\subsection{Batdongsan.com.vn}

Batdongsan.com.vn~\cite{Batdongsan} là một trong những nền tảng rao vặt bất động sản lớn và lâu đời nhất tại Việt Nam, cung cấp tin đăng cho cả nhu cầu mua bán lẫn cho thuê trên phạm vi toàn quốc. Điểm mạnh của nền tảng là kho tin đăng lớn, công cụ tìm kiếm và lọc theo nhiều tiêu chí (khu vực, khoảng giá, diện tích, loại hình), cùng mô hình kinh doanh dựa trên các gói tin VIP và đẩy tin giúp người đăng tăng khả năng tiếp cận. Tuy nhiên, nền tảng thiên về phục vụ môi giới chuyên nghiệp và phân khúc mua bán; trải nghiệm cho người thuê ở phân khúc nhà trọ, phòng trọ giá rẻ chưa phải là trọng tâm, và chất lượng tin đăng phụ thuộc nhiều vào người đăng.

\FloatBarrier
\subsection{Phongtro123}

Phongtro123~\cite{Phongtro123} là nền tảng tập trung chuyên biệt vào phân khúc \textit{phòng trọ, nhà trọ} và các loại hình cho thuê giá rẻ, hướng tới đối tượng sinh viên và người lao động. Giao diện đơn giản, dễ sử dụng, cho phép lọc nhanh theo khu vực và mức giá, phù hợp với nhu cầu tìm trọ phổ thông. Bù lại, nền tảng chủ yếu dừng ở vai trò bảng tin rao vặt: thông tin tin đăng còn sơ sài, thiếu các công cụ hỗ trợ ra quyết định như so sánh, định giá tham khảo, bản đồ trực quan hay gợi ý cá nhân hóa, và việc kiểm soát chất lượng/độ tin cậy của tin đăng còn hạn chế.

\FloatBarrier
\subsection{Nhà Tốt}

Nhà Tốt~\cite{NhaTot} (chuyên mục bất động sản của nền tảng rao vặt Chợ Tốt) cung cấp tin đăng cho thuê và mua bán theo mô hình rao vặt giữa người dùng với người dùng. Nhờ thừa hưởng lượng người dùng lớn của Chợ Tốt, nền tảng có lưu lượng truy cập cao và đa dạng loại hình. Nền tảng có hỗ trợ bản đồ và một số tiện ích tìm kiếm, song giống các nền tảng nội địa khác, mức độ ứng dụng trí tuệ nhân tạo để định giá, gợi ý hay hỗ trợ hội thoại vẫn còn mờ nhạt, và trải nghiệm còn nặng tính rao vặt thủ công.

\FloatBarrier
\subsection{Zillow}

Zillow~\cite{Zillow} là nền tảng bất động sản hàng đầu tại Hoa Kỳ, nổi bật với việc ứng dụng dữ liệu và công nghệ ở mức cao. Tính năng tiêu biểu là \textit{Zestimate}, mô hình ước tính giá trị bất động sản tự động dựa trên dữ liệu thị trường, cùng với bản đồ giàu thông tin, dữ liệu khu vực chi tiết và trải nghiệm người dùng được trau chuốt. Zillow cho thấy tiềm năng to lớn của việc tích hợp dữ liệu và AI vào lĩnh vực bất động sản. Tuy vậy, nền tảng được thiết kế cho thị trường Hoa Kỳ với đặc thù dữ liệu, pháp lý và hành vi người dùng khác biệt, do đó không trực tiếp phù hợp với thị trường thuê trọ tại Việt Nam.

\FloatBarrier
\subsection{So sánh ưu nhược điểm}

Bảng~\ref{tab:compare-platforms} tổng hợp một số tiêu chí so sánh giữa các nền tảng đã khảo sát.

\begin{table}[!ht]
\centering
\small
\caption{So sánh tính năng giữa đề tài và các nền tảng hiện có}
\label{tab:compare-platforms}
\begin{tabular}{|L{0.22\textwidth}|C{0.11\textwidth}|C{0.11\textwidth}|C{0.10\textwidth}|C{0.10\textwidth}|C{0.13\textwidth}|}
\hline
\textbf{Tiêu chí} & \textbf{Batdong\-san} & \textbf{Phong\-tro123} & \textbf{Nhà Tốt} & \textbf{Zillow} & \textbf{Thuê Nhà Trọ} \\
\hline
Tập trung nhà trọ giá rẻ & Trung bình & Cao & Trung bình & Thấp & Cao \\
\hline
Tìm kiếm \& lọc đa tiêu chí & Cao & Trung bình & Trung bình & Cao & Cao \\
\hline
Bản đồ trực quan & Có & Hạn chế & Có & Cao & Có \\
\hline
Định giá thông minh & Không & Không & Không & Có & Có \\
\hline
Gợi ý cá nhân hóa & Hạn chế & Không & Hạn chế & Có & Có \\
\hline
Trợ lý hội thoại thông minh & Không & Không & Không & Hạn chế & Có \\
\hline
Phù hợp thị trường Việt Nam & Cao & Cao & Cao & Thấp & Cao \\
\hline
\end{tabular}
\end{table}

Qua so sánh có thể thấy: các nền tảng nội địa (Batdongsan.com.vn, Phongtro123, Nhà Tốt) phù hợp với thị trường và thói quen người dùng Việt Nam nhưng hầu như chưa khai thác trí tuệ nhân tạo để nâng cao trải nghiệm tìm kiếm và ra quyết định; trong khi đó, nền tảng quốc tế Zillow thể hiện rõ giá trị của việc tích hợp dữ liệu và AI nhưng lại không được thiết kế cho bối cảnh Việt Nam. Thuê Nhà Trọ lấp đầy khoảng trống đó: đặt trọng tâm vào phân khúc nhà trọ giá rẻ tại Việt Nam, đồng thời tích hợp đầy đủ các năng lực AI (định giá tham khảo, gợi ý cá nhân hóa và trợ lý hội thoại) mà các nền tảng hiện có chưa đáp ứng được.

\section{Cơ sở lý thuyết về trí tuệ nhân tạo}

Để khắc phục các hạn chế của phương pháp truyền thống, đề tài nghiên cứu và ứng dụng một số công nghệ trí tuệ nhân tạo, bao gồm các mô hình ngôn ngữ lớn, kỹ thuật tạo sinh kết hợp truy xuất, tác nhân AI và các mô hình gợi ý cá nhân hóa.

\FloatBarrier
\subsection{Mô hình ngôn ngữ lớn và Kỹ thuật Prompting}

Mô hình ngôn ngữ lớn là các mạng nơ-ron sâu dựa trên kiến trúc Transformer~\cite{Vaswani2017}, được huấn luyện trên các tập dữ liệu văn bản khổng lồ để hiểu và tạo sinh ngôn ngữ tự nhiên. Trong đề tài này, mô hình Gemini của Google được sử dụng làm nền tảng trí tuệ cốt lõi nhờ khả năng xử lý đa phương tiện và hiệu năng suy luận tốt \cite{GeminiDocs}.

Trong quá trình tích hợp LLM, kỹ thuật thiết kế chỉ thị đóng vai trò quyết định để định hình hành vi phản hồi của mô hình. Nhóm nghiên cứu đã ứng dụng hai kỹ thuật chính:
\begin{itemize}
    \item \textbf{In-Context Learning (Học ngữ cảnh / Few-Shot Prompting):} Thay vì tinh chỉnh lại trọng số mô hình vốn tốn kém tài nguyên, hệ thống cung cấp trực tiếp các ví dụ mẫu (inputs-outputs) hoặc các dữ liệu thực tế lân cận ngay trong prompt. Điều này giúp mô hình nhận diện cấu trúc, phong cách văn phong và quy luật định giá để phản hồi chính xác~\cite{Brown2020}.
    \item \textbf{System Prompting:} Thiết lập một tập quy tắc hành xử cố định cho mô hình.
\end{itemize}

\FloatBarrier
\subsection{Tạo sinh kết hợp truy xuất}

Mô hình ngôn ngữ lớn tuy rất thông minh nhưng lại mắc phải hai hạn chế nghiêm trọng: tri thức tĩnh (chỉ giới hạn trong thời điểm dữ liệu huấn luyện) và hiện tượng ``ảo tưởng'' (tự tạo ra thông tin sai lệch). Nhằm khắc phục hạn chế này, kỹ thuật tạo sinh kết hợp truy xuất đã được phát triển~\cite{Lewis2020}.

Quy trình RAG hoạt động theo ba bước chính:
\begin{enumerate}
    \item \textbf{Truy xuất (Retrieve):} Khi người dùng đặt câu hỏi (ví dụ: ``thủ tục đặt cọc nhà như thế nào?''), hệ thống sẽ chuyển hóa câu hỏi này và thực hiện tìm kiếm ngữ nghĩa trên cơ sở dữ liệu tri thức có sẵn (được lưu dưới dạng các tài liệu luật cư trú, cẩm nang thuê phòng trọ).
    \item \textbf{Tăng cường (Augment):} Đoạn văn bản kết quả tìm kiếm có độ liên quan cao nhất sẽ được đính kèm vào prompt gửi cho LLM như một tài liệu tham khảo chính thức.
    \item \textbf{Tạo sinh (Generate):} LLM đọc tài liệu này để tổng hợp câu trả lời chính xác, giúp hạn chế hiện tượng đoán mò và bảo đảm câu trả lời bám sát quy định pháp luật hoặc thông tin thực tế của hệ thống.
\end{enumerate}

\FloatBarrier
\subsection{Tác nhân trí tuệ nhân tạo và Vòng lặp ReAct}

Tác nhân AI là một thực thể phần mềm thông minh có khả năng tự động hóa quy trình ra quyết định để thực hiện một mục tiêu cụ thể bằng cách quan sát môi trường và kích hoạt các công cụ hành động (tools) có sẵn. Trong đề tài, trợ lý chatbot tìm phòng được xây dựng dưới dạng một Agent hoạt động theo vòng lặp ReAct (Reasoning and Acting)~\cite{Yao2023ReAct}.

Vòng lặp ReAct kết hợp đồng thời khả năng lập luận (Reasoning) và hành động (Acting) của LLM qua bốn giai đoạn lặp đi lặp lại:
\begin{itemize}
    \item \textbf{Suy nghĩ (Thought):} Agent phân tích câu lệnh của người dùng để xác định bước tiếp theo.
    \item \textbf{Hành động (Action):} Agent quyết định gọi một công cụ cụ thể với tham số đầu vào được trích xuất từ suy nghĩ.
    \item \textbf{Quan sát (Observation):} Hệ thống thực thi công cụ trong cơ sở dữ liệu và trả về kết quả thô cho Agent quan sát.
    \item \textbf{Lặp lại / Trả lời (Response):} Agent đánh giá kết quả quan sát. Nếu đã đủ thông tin, Agent sẽ sinh câu trả lời tự nhiên để phản hồi cho người dùng; nếu chưa đủ, Agent tiếp tục lập luận để thực hiện một hành động khác. Để tránh vòng lặp vô hạn do lỗi suy luận, hệ thống cấu hình giới hạn tối đa 12 bước thực thi cho mỗi câu hỏi.
\end{itemize}

\FloatBarrier
\subsection{Hệ thống gợi ý lai}

Hệ thống gợi ý tin đăng giúp cá nhân hóa trang chủ của từng người dùng dựa trên sở thích và hành vi của họ. Hệ thống Thuê Nhà Trọ triển khai cơ chế gợi ý lai~\cite{Burke2002} kết hợp hai hướng tiếp cận chính:
\begin{itemize}
    \item \textbf{Lọc dựa trên nội dung:} Gợi ý các tin đăng có đặc trưng tương đồng (tiện ích, khoảng giá, vị trí) với các tin đăng mà người dùng đó đã từng xem, lưu hoặc tương tác nhiều trong lịch sử. Độ tương đồng giữa các đặc trưng được tính toán bằng chỉ số độ tương đồng Cosine trên không gian vector.
    \item \textbf{Lọc cộng tác:} Tìm kiếm các người dùng có hành vi tương tự (lịch sử xem tin giống nhau) để gợi ý chéo các tin đăng mà người dùng tương tự kia đã thích nhưng người dùng hiện hành chưa biết tới.
\end{itemize}
Hai thuật toán này được tích hợp song song và chuẩn hóa điểm số, kết hợp với các hệ số hiệu chỉnh đặc thù (như khoảng cách địa lý từ cơ sở làm việc/học tập của người dùng và mức độ ưu tiên đẩy tin của tài khoản VIP) để xếp hạng danh sách gợi ý cuối cùng.

\section{Khoảng trống và cơ hội ứng dụng AI trong lĩnh vực cho thuê nhà trọ}

Qua khảo sát hiện trạng tại các nền tảng cho thuê bất động sản hiện nay (Batdongsan, Phongtro123, Nhà Tốt), nhóm nghiên cứu nhận diện rõ những khoảng trống công nghệ lớn và cơ hội ứng dụng trí tuệ nhân tạo để nâng tầm trải nghiệm của thị trường:

\begin{itemize}
    \item \textbf{Thay thế bộ lọc cứng nhắc bằng Giao diện hội thoại tự nhiên:} Các hệ thống hiện tại bắt buộc người thuê phải chọn hàng loạt bộ lọc phức tạp và thường dẫn tới kết quả rỗng khi điều kiện lọc quá khắt khe. Ứng dụng AI Agent cho phép người dùng mô tả nhu cầu dưới dạng một câu nói tự nhiên, Agent có khả năng tự động nới lỏng hoặc đề xuất phương án thay thế thông minh (ví dụ: gợi ý phòng ở quận lân cận khi quận đích đã hết phòng), mang lại trải nghiệm tìm kiếm mượt mà như đang giao tiếp trực tiếp với môi giới thật.
    \item \textbf{Giải quyết bài toán định giá thiếu minh bạch:} Người thuê và người cho thuê hiện nay định giá bất động sản chủ yếu dựa trên cảm tính hoặc tốn nhiều thời gian tự dò tìm so sánh thủ công. Tích hợp AI định giá kết hợp In-Context Learning giúp tự động phân tích hàng ngàn dữ liệu so sánh lân cận thời gian thực để đưa ra một báo cáo ước tính giá khách quan tức thời, hạn chế tối đa việc ép giá hoặc định giá quá cao.
    \item \textbf{Tự động hóa hoàn toàn quy trình kiểm duyệt chất lượng:} Quy trình kiểm duyệt tin thủ công bằng con người của các sàn truyền thống thường gây trễ bài đăng (từ vài giờ đến vài ngày) và dễ bỏ sót tin spam hoặc hình ảnh không phù hợp. Ứng dụng mô hình Gemini Multimodal cho phép quét lập tức hình ảnh, video và văn bản ngay khi chủ trọ vừa bấm đăng tin để phát hiện tin trùng, ảnh mờ, hoặc thông tin rác, giúp tăng tốc độ phê duyệt tin đăng lên thời gian thực trong khi vẫn đảm bảo độ tin cậy của dữ liệu trên sàn.
    \item \textbf{Cá nhân hóa tối đa trải nghiệm người dùng:} Phản lớn các trang rao vặt hiển thị tin đăng theo thứ tự thời gian hoặc mức phí dịch vụ thuần túy, dẫn tới việc người dùng bị ngợp bởi thông tin không liên quan. Việc triển khai hệ thống gợi ý lai (Collaborative + Content-based) giúp cá nhân hóa giao diện trang chủ theo đúng nhu cầu cư trú riêng biệt của từng cá nhân (khoảng cách đi lại tối ưu, sở thích tiện nghi), cải thiện tỷ lệ chuyển đổi giao dịch cho nền tảng.
\end{itemize}
